\section{Background}

The proliferation of autonomous multi-agent systems across safety-critical domains---including healthcare~\cite{bansal2019}, autonomous vehicles, and financial infrastructure---has catalyzed renewed interest in accountability mechanisms that can operate at machine speed without sacrificing human oversight. As agentic systems evolve from isolated, single-purpose tools into interconnected ecosystems capable of hierarchical reasoning and independent action selection, the design of fail-safe architectures and human-in-the-loop (HITL) safeguards becomes both more urgent and more challenging. This section establishes the theoretical and historical foundations upon which the Bailout Protocol is constructed, reviewing accountability frameworks in multi-agent AI, HITL design principles, and established fail-safe design traditions, before introducing the BX3 Framework's three-layer model as an organizing architecture for agentic governance.

% -----------------------------------------------------------------------------
\subsection{Accountability in Multi-Agent AI Systems}

Classical accountability theory, rooted in organizational sociology and legal philosophy, distinguishes between \emph{forward-looking} accountability (incentivizing future compliance) and \emph{backward-looking} accountability (assigning responsibility after harm)~\cite{wiener1948}. In multi-agent systems, both dimensions present unique challenges. Agents may operate as distributed processes where no single node possesses complete situational awareness, and harm may emerge from emergent interactions among otherwise individually benign components.

Dijkstra's foundational work on cooperative sequential processes~\cite{dijkstra1982} established that reasoning about composed systems requires formal verification boundaries. When agents share resources, communicate via asynchronous channels, or operate under partial information, the composition itself becomes a source of non-determinism that cannot be attributed to any single agent. This compositional opacity is a direct threat to accountability: if neither the designer nor the operator can trace a causal chain from system behavior to individual agent decisions, backward-looking responsibility assignment becomes intractable.

Modern multi-agent accountability frameworks~\cite{bansal2019} address this through structured audit trails, decision-point logging, and causal inference over agent histories. However, these approaches assume that agents are designed with accountability in mind \emph{ab initio}---a property not guaranteed in loosely integrated systems or third-party components. The Bailout Protocol is motivated by the observation that \emph{restoration of accountability after failure} is as important as prevention: when an agent or agent subsystem enters an unanticipated state, the system must be capable of returning to a configuration where human oversight can be re-established, and where the causal history of the failure is preserved for retrospective analysis.

% -----------------------------------------------------------------------------
\subsection{Human-in-the-Loop Design}

Human-in-the-loop architectures have long been recognized as essential for safety-critical systems. Norman~\cite{wiener1948} articulated the danger of ceding consequential decisions to autonomous processes without preserving meaningful human agency---a concern that grows more acute as AI systems approach and exceed human-level performance in narrow domains. Leveson's work on safety engineering formalizes this into a principle: systems must maintain the ability to detect and respond to novel hazards, and this responsiveness must ultimately reside with humans who can exercise judgment beyond the envelope of the system's verified operating conditions.

In the context of agentic AI, HITL design encounters a structural tension between two requirements: (\emph{i}) the need for low-latency autonomous action to prevent harm in time-critical scenarios, and (\emph{ii}) the need for human authorization to prevent irreversible or catastrophic outcomes. This tension is sometimes framed as the \emph{speed--safety tradeoff}, but more recent treatments~\cite{bansal2019} reject the tradeoff as a false dichotomy. Rather, they propose layered architectures where HITL is enforced not uniformly across all decisions but as a function of consequence severity and decision reversibility.

Consequence-severity stratification allows agents to operate autonomously within bounded risk envelopes while reserving human review for decisions that exceed defined thresholds. Reversibility considerations further refine this: actions that are easily undone (e.g., data queries, content recommendations) can be automated with lower oversight, while actions with irreversible downstream effects (e.g., financial transactions, system configuration changes, physical actuation) require explicit human authorization or confirmation.

The design implication is that agentic systems must implement \emph{proportional oversight}: the depth and form of human involvement scales with the consequence profile of each action class. Static, binary override switches (``kill switches'') address the most extreme cases but do not support this proportional model. The Bailout Protocol's revocation mechanism is designed to operationalize proportional oversight by providing a graded response surface---from suspension of specific capabilities to full agent shutdown---that can be triggered by human overseers in real time.

% -----------------------------------------------------------------------------
\subsection{Fail-Safe Design in Safety-Critical Systems}

The concept of fail-safe design originates in control systems engineering and has been codified in numerous safety standards, including IEC 61508 and its domain-specific derivatives. A fail-safe system is one that, upon experiencing a defined failure mode, transitions to a state that minimizes harm to persons, equipment, or the environment. Classical fail-safe design often leverages physical redundancy (e.g., backup sensors, watchdog timers) and invariant enforcement (e.g., dead man's switches, limit checkers).

Trist's analysis of organizational resilience~\cite{tristr81} extends fail-safe thinking from physical systems to sociotechnical systems, arguing that resilient organizations are those that preserve core functions under stress rather than optimizing exclusively for steady-state performance. Applied to agentic systems, this suggests that fail-safe design must do more than detect and respond to failures; it must ensure that the system's \emph{governance infrastructure}---the mechanisms by which humans direct, constrain, and override agents---remains operational under stress.

Wiener's cybernetics framework~\cite{wiener1948} provides a unifying lens: feedback loops are the fundamental mechanism through which systems maintain stability in the presence of disturbance. For agentic AI, the human-in-the-loop constitutes a feedback channel from the operator's world-model to the agent's action selection process. A fail-safe agentic system must therefore treat the HITL channel as a \emph{critical dependency} that must be protected, monitored, and rapidly restorable.

This perspective motivates the Bailout Protocol's emphasis on \emph{detecting degradation of the oversight channel}, not just failures within the agent subsystem. An agent that continues to operate but whose human oversight channel has been compromised represents a more fundamental safety failure than an agent that simply stops. The Protocol's heartbeat mechanism and idempotent revocation procedures are concrete implementations of this principle.

% -----------------------------------------------------------------------------
\subsection{The BX3 Framework: A Three-Layer Architecture for Agentic Governance}

The BX3 Framework organizes agentic system governance into three conceptual layers, each with distinct responsibilities and failure modes:

\begin{description}
\item[\textbf{Layer 1 -- Directive:}] The strategic layer, encoding goals, policies, and ethical boundaries as explicit, machine-evaluable constraints. This layer defines the ``north star'' of agent behavior and is the primary locus of human authority. Failures at this layer indicate misalignment between operator intent and agent world-model.
\item[\textbf{Layer 2 -- Execution:}] The tactical layer, responsible for decomposing directives into actionable plans, selecting among candidate actions, and managing cross-agent coordination. This layer translates intent into operation and is the primary locus of autonomous decision-making. Failures at this layer often manifest as policy violations, resource conflicts, or decision-quality degradation.
\item[\textbf{Layer 3 -- Substrate:}] The operational layer, encompassing the computational infrastructure, network communications, tool access, and environment interfaces on which execution operates. This layer is the foundation upon which all higher-layer functions depend. Failures at this layer include infrastructure outages, tool malfunctions, and data integrity violations.
\end{description}

The three-layer model enables \emph{failure localization}: when an adverse event occurs, the framework constrains the search space of causal explanations to the layer in which the failure originates, while also mandating verification that co-located layers have not been compromised by propagation. This architectural discipline supports the construction of audit trails that are meaningful for retrospective accountability and enables the graded intervention model that the Bailout Protocol implements.

Critically, the BX3 three-layer model does not treat these layers as independent silos but as a tightly coupled hierarchy where each layer provides services to the layer above and receives guarantees from the layer below. The Directive layer relies on the Execution layer to faithfully interpret and implement its constraints; the Execution layer relies on the Substrate to provide reliable tool access and environment feedback. Fail-safe design at each layer requires that this reliance be bounded and monitored.

The Bailout Protocol extends the BX3 Framework by providing a cross-layer failure response protocol that coordinates revocation of agent capabilities across all three layers simultaneously, ensuring that when a failure is detected at any layer, the entire agentic system transitions to a safe, human-oversight-resumable state rather than continuing to operate in a degraded or ambiguous mode.